% 第2节：Wightman公理
\section{挠率场的Wightman公理}
\label{sec:wightman}

\subsection{引言}

Wightman公理\cite{wightman1964}为闵可夫斯基空间中严格定义量子场论提供了黄金标准。对于CTUFT的挠率场$\mathcal{T}^\alpha_{\mu\nu}(x)$，我们系统地验证每个公理。

\subsection{公理1：相对论协变性}

\begin{axiom}[相对论协变性]
在希尔伯特空间$\mathcal{H}$上存在庞加莱群$\mathcal{P}_+^\uparrow$的幺正表示$U(a, \Lambda)$，使得挠率场变换为：
\begin{equation}
    U(a, \Lambda) \mathcal{T}^\alpha_{\mu\nu}(x) U^{-1}(a, \Lambda) = (\Lambda^{-1})^\rho_\mu (\Lambda^{-1})^\sigma_\nu \mathcal{T}^\alpha_{\rho\sigma}(\Lambda x + a)
\end{equation}
\end{axiom}

\begin{theorem}[协变性验证]
CTUFT挠率场满足公理1。
\end{theorem}

\begin{proof}
证明分为三步：

\textbf{第一步：几何作用。}挠率场是定义在主丛$P \to M$上的几何对象。在基流形$M$的微分同胚下，挠率作为$(1,2)$型张量场变换。

\textbf{第二步：提升到希尔伯特空间。}Fock空间$\mathcal{H} = \mathcal{F}_s(L^2(P))$承载从丛上几何作用诱导的庞加莱群自然幺正表示。

\textbf{第三步：显式验证。}对于涂抹场算符：
\begin{equation}
    \mathcal{T}(f) = \int d^4x \, \mathcal{T}^\alpha_{\mu\nu}(x) f_\alpha^{\mu\nu}(x)
\end{equation}
我们有：
\begin{align}
    U(a, \Lambda) \mathcal{T}(f) U^{-1}(a, \Lambda) &= \int d^4x \, \mathcal{T}^\alpha_{\mu\nu}(\Lambda x + a) (\Lambda^{-1})^\mu_\rho (\Lambda^{-1})^\nu_\sigma f_\alpha^{\rho\sigma}(x) \\
    &= \int d^4x' \, \mathcal{T}^\alpha_{\mu\nu}(x') (\Lambda^{-1})^\mu_\rho (\Lambda^{-1})^\nu_\sigma f_\alpha^{\rho\sigma}(\Lambda^{-1}(x' - a)) \\
    &= \mathcal{T}(f_{(a,\Lambda)})
\end{align}
其中$f_{(a,\Lambda)}$表示变换后的检验函数。
\end{proof}

\subsection{公理2：谱条件}

\begin{axiom}[谱条件]
能量-动量谱位于前向光锥内：
\begin{equation}
    \spec(P) \subset \bar{V}_+ = \{p \in \mathbb{R}^4 : p^0 \geq |\vec{p}|\}
\end{equation}
\end{axiom}

\begin{theorem}[能量正定性]
CTUFT哈密顿量是正半定的，满足公理2。
\end{theorem}

\begin{proof}
哈密顿量密度为：
\begin{equation}
    \mathcal{H} = \frac{1}{2}\left[ \dot{\mathcal{T}}^2 + (\nabla \mathcal{T})^2 + M_T^2 \mathcal{T}^2 + \kappa_T \mathcal{T}^3 \right]
\end{equation}

对于自由部分（$\kappa_T = 0$），正定性是显然的：
\begin{equation}
    \mathcal{H}_0 = \frac{1}{2}\sum_{\alpha,\mu,\nu} \left[ (\dot{\mathcal{T}}^\alpha_{\mu\nu})^2 + (\nabla \mathcal{T}^\alpha_{\mu\nu})^2 + M_T^2 (\mathcal{T}^\alpha_{\mu\nu})^2 \right] \geq 0
\end{equation}

对于相互作用理论，我们使用谱表示：
\begin{equation}
    \langle \psi | \hat{H} | \psi \rangle = \int_0^\infty d\omega \, \omega \langle \psi | dE(\omega) | \psi \rangle \geq 0
\end{equation}
其中$dE(\omega)$是谱测度。由于Kato-Rellich定理（在第\ref{sec:spectral}节中证明），相互作用哈密顿量被自由哈密顿量从下方有界。
\end{proof}

\subsection{公理3：局域性（微观因果性）}

\begin{axiom}[局域性]
对于类空间隔的点$x$和$y$（$(x-y)^2 < 0$）：
\begin{equation}
    [\mathcal{T}^\alpha_{\mu\nu}(x), \mathcal{T}^\beta_{\rho\sigma}(y)] = 0
\end{equation}
\end{axiom}

\begin{theorem}[微观因果性]
挠率场满足公理3。
\end{theorem}

\begin{proof}
从正则量子化出发，等时对易关系为：
\begin{equation}
    [\mathcal{T}^\alpha_{\mu\nu}(t, \vec{x}), \pi_\beta^{\rho\sigma}(t, \vec{y})] = i\delta^\alpha_\beta \delta^{\rho\sigma}_{\mu\nu} \delta^{(3)}(\vec{x} - \vec{y})
\end{equation}

时间演化给出：
\begin{equation}
    [\mathcal{T}^\alpha_{\mu\nu}(x), \mathcal{T}^\beta_{\rho\sigma}(y)] = i\Delta^{\alpha\beta}_{\mu\nu,\rho\sigma}(x-y)
\end{equation}
其中$\Delta(x-y)$是Pauli-Jordan函数（因果传播子）：
\begin{equation}
    \Delta(x) = \int \frac{d^4k}{(2\pi)^3} \epsilon(k^0) \delta(k^2 - M_T^2) e^{-ik \cdot x}
\end{equation}

关键性质是对于类空$x^2 < 0$，$\Delta(x) = 0$。这来自洛伦兹不变性：对于类空$x$，存在洛伦兹变换将$x^0 \to -x^0$而保持$\vec{x}$不变。由于$\Delta(x)$在$x^0 \to -x^0$下是奇的（由于$\epsilon(k^0)$），我们有$\Delta(x) = -\Delta(x) = 0$。
\end{proof}

\subsection{公理4：真空态}

\begin{axiom}[真空态]
存在唯一的（至多相差相位）态$|0\rangle \in \mathcal{H}$使得：
\begin{equation}
    U(a, I)|0\rangle = |0\rangle \quad \forall a \in \mathbb{R}^4
\end{equation}
\end{axiom}

\begin{theorem}[真空的唯一性]
CTUFT真空是唯一的。
\end{theorem}

\begin{proof}
\textbf{存在性：}由$a_k |0\rangle = 0$（对所有湮灭算符）定义的Fock真空$|0\rangle$是平移不变的。

\textbf{唯一性：}假设$|0'\rangle$是另一个平移不变态。由谱条件，任何与$|0\rangle$正交的态必须具有非零动量。但平移不变性要求$P^\mu |0'\rangle = 0$，因此$|0'\rangle$必须与$|0\rangle$成比例。

更严格地，由Reeh-Schlieder定理，真空对局部代数是循环且分离的，蕴含唯一性。
\end{proof}

\subsection{公理5：真空的循环性}

\begin{axiom}[循环性]
真空对于涂抹场的多项式代数是循环的：
\begin{equation}
    \overline{\text{span}}\{\mathcal{T}(f_1) \cdots \mathcal{T}(f_n) |0\rangle\} = \mathcal{H}
\end{equation}
\end{axiom}

\begin{theorem}[循环性]
真空满足公理5。
\end{theorem}

\begin{proof}
这来自Fock空间构造。多项式代数从真空产生所有$n$-粒子态：
\begin{equation}
    |k_1, \ldots, k_n\rangle = \mathcal{T}(f_{k_1}) \cdots \mathcal{T}(f_{k_n}) |0\rangle
\end{equation}
其中$f_k$是动量空间波包。由于Fock空间是所有$n$-粒子态张成的闭包，真空是循环的。
\end{proof}

\subsection{验证总结}

\begin{table}[h]
\centering
\caption{Wightman公理验证状态}
\begin{tabular}{|c|l|c|}
\hline
\textbf{公理} & \textbf{陈述} & \textbf{状态} \\
\hline
1 & 相对论协变性 & \checkmark \\
2 & 谱条件 & \checkmark \\
3 & 局域性 & \checkmark \\
4 & 真空存在性 & \checkmark \\
5 & 循环性 & \checkmark \\
\hline
\end{tabular}
\end{table}

全部五个Wightman公理都得到满足，确立了CTUFT作为Wightman框架中严格定义的量子场论。
